\section{Energ\'ia de horizonte y celda de media acci\'on}
\label{rm:chap:horizonte-media-accion}

\subsection{Convenciones geom\'etricas del parche de Sitter}

Consideremos un espacio de Sitter cuadridimensional con radio areal
\(R>0\), velocidad causal \(c\) y constante gravitatoria \(G\). Adoptamos
la convenci\'on
\begin{equation}
 H_R=\frac cR,
 \qquad
 \Lambda_R=\frac3{R^2},
 \qquad
 \ell_P^2=\frac{G\hbar}{c^3}.
 \label{rm:eq:horizon-ds-conventions}
\end{equation}
La densidad de energ\'ia asociada al t\'ermino cosmol\'ogico es
\begin{equation}
 \varepsilon_\Lambda
 :=\frac{\Lambda_Rc^4}{8\pi G}.
 \label{rm:eq:horizon-vacuum-density}
\end{equation}
La secci\'on espacial plana del parche contiene, hasta el radio \(R\), el
volumen
\begin{equation}
 V_R=\frac{4\pi}{3}R^3.
 \label{rm:eq:horizon-flat-volume}
\end{equation}
Estas convenciones fijan positivamente la energ\'ia cuasilocal usada en el
cap\'itulo. La elecci\'on de otro signo para la carga de Killing de Sitter
pertenece a una construcci\'on distinta y debe traducirse antes de comparar
f\'ormulas.

\subsection{Energ\'ia de vac\'io y energ\'ia de Misner--Sharp}

En una geometr\'ia esf\'ericamente sim\'etrica, sea
\begin{equation}
 \dd s^2=g_{ab}(x)\dd x^a\dd x^b
 +\mathcal R(x)^2\dd\Omega_2^2.
 \label{rm:eq:horizon-spherical-metric}
\end{equation}
La energ\'ia de Misner--Sharp \cite{rm:MisnerSharp1964}, en la convenci\'on positiva adoptada, se
define por
\begin{equation}
 E_{\rm MS}(\mathcal R)
 :=\frac{c^4\mathcal R}{2G}
 \left(1-g^{ab}\nabla_a\mathcal R\nabla_b\mathcal R\right).
 \label{rm:eq:horizon-Misner-Sharp}
\end{equation}
Una esfera marginal satisface
\begin{equation}
 g^{ab}\nabla_a\mathcal R\nabla_b\mathcal R=0.
 \label{rm:eq:horizon-marginal-sphere}
\end{equation}

\begin{theorem}[Confluencia volum\'etrica y cuasilocal]
\label{rm:thm:horizon-vacuum-MS}
En el horizonte de Sitter \(\mathcal R=R\), la energ\'ia integrada del
vac\'io y la energ\'ia de Misner--Sharp coinciden exactamente:
\begin{equation}
 \boxed{
 E_\Lambda
 :=\varepsilon_\Lambda V_R
 =E_{\rm MS}(R)
 =\frac{c^4R}{2G}.}
 \label{rm:eq:horizon-E-Lambda}
\end{equation}
\end{theorem}

\begin{proof}
Las ecuaciones
\cref{rm:eq:horizon-ds-conventions,rm:eq:horizon-vacuum-density,rm:eq:horizon-flat-volume}
dan
\[
 \varepsilon_\Lambda V_R
 =\frac{3c^4}{8\pi GR^2}\frac{4\pi R^3}{3}
 =\frac{c^4R}{2G}.
\]
La condici\'on marginal \cref{rm:eq:horizon-marginal-sphere} aplicada a
\cref{rm:eq:horizon-Misner-Sharp} produce el mismo resultado.
\end{proof}

La igualdad re\'une dos publicaciones geom\'etricas. La primera integra una
densidad volum\'etrica; la segunda eval\'ua el defecto de atrapamiento de la
esfera areal. Ambas conservan el radio \(R\) y la compliancia gravitatoria
\(G/c^4\).

\subsection{Temperatura, entrop\'ia y producto de horizonte}

La gravedad superficial de Sitter tiene m\'odulo \(c^2/R\). La temperatura
de Gibbons--Hawking \cite{rm:GibbonsHawking1977} y la entrop\'ia de
Bekenstein--Hawking \cite{rm:Bekenstein1973,rm:Hawking1975} son
\begin{align}
 k_BT_H&=\frac{\hbar c}{2\pi R},
 \label{rm:eq:horizon-temperature}\\
 \frac{S_H}{k_B}
 &=\frac{A_H}{4\ell_P^2}
 =\frac{4\pi R^2}{4\ell_P^2}
 =\frac{\pi R^2}{\ell_P^2}.
 \label{rm:eq:horizon-entropy}
\end{align}

\begin{theorem}[Identidad termodin\'amica de Sitter]
\label{rm:thm:horizon-TS}
Para todo radio \(R\) de la familia declarada,
\begin{equation}
 \boxed{
 T_HS_H=\frac{c^4R}{2G}=E_\Lambda.}
 \label{rm:eq:horizon-TS-E}
\end{equation}
\end{theorem}

\begin{proof}
Multiplicar \cref{rm:eq:horizon-temperature,rm:eq:horizon-entropy} da
\[
 T_HS_H
 =\frac{\hbar c}{2\pi Rk_B}
  \frac{k_B\pi R^2}{\ell_P^2}
 =\frac{\hbar cR}{2\ell_P^2}.
\]
La identidad \(\ell_P^2=G\hbar/c^3\) convierte el \'ultimo miembro en
\(c^4R/(2G)\), igual a \cref{rm:eq:horizon-E-Lambda}.
\end{proof}

El factor \(1/2\) posee una descomposici\'on geom\'etrica transparente. El
factor areal \(4\pi\), la normalizaci\'on hologr\'afica por cuatro y la
periodicidad t\'ermica \((2\pi)^{-1}\) satisfacen
\begin{equation}
 \frac{4\pi}{4}\frac1{2\pi}=\frac12.
 \label{rm:eq:horizon-half-factor}
\end{equation}
La mitad es, por tanto, la contracci\'on conjunta entre cuantificaci\'on
areal y periodicidad eucl\'idea.

\subsection{Cuanto completo del horizonte y fuerza de Planck}

Definimos la tensi\'on o fuerza de Planck y el cuanto completo asociado al
radio por
\begin{equation}
 \mathscr F_P:=\frac{c^4}{G},
 \qquad
 \boxed{
 \mathcal Q_H(R):=2T_HS_H
 =\frac{c^4R}{G}=\mathscr F_P\,R.}
 \label{rm:eq:horizon-full-quantum}
\end{equation}
La energ\'ia cuasilocal ocupa media celda del cuanto completo:
\begin{equation}
 \boxed{E_\Lambda=T_HS_H=\frac{\mathcal Q_H(R)}2.}
 \label{rm:eq:horizon-half-quantum-general}
\end{equation}

\begin{proposition}[Compliancia radial]
\label{rm:prop:horizon-compliance}
La derivada del cuanto completo respecto del radio es la fuerza de Planck,
y la derivada de la energ\'ia cuasilocal es su mitad:
\begin{equation}
 \frac{\dd\mathcal Q_H}{\dd R}=\mathscr F_P,
 \qquad
 \frac{\dd E_\Lambda}{\dd R}=\frac{\mathscr F_P}{2}.
 \label{rm:eq:horizon-radial-derivatives}
\end{equation}
La compliancia gravitatoria \(\mathscr C_P=G/c^4\) satisface
\begin{equation}
 R=\mathscr C_P\mathcal Q_H.
 \label{rm:eq:horizon-compliance-law}
\end{equation}
\end{proposition}

\begin{proof}
Las tres identidades son derivadas o inversiones directas de
\cref{rm:eq:horizon-full-quantum}.
\end{proof}

La ecuaci\'on \cref{rm:eq:horizon-compliance-law} expresa la respuesta radial:
el cuanto completo aplica tensi\'on y la geometr\'ia publica un radio mediante
la compliancia universal. El producto \(T_HS_H\) ocupa la mitad cuasilocal
de esa respuesta.

\subsection{Especializaci\'on al cierre cronol\'ogico de 108 estados}

El cierre cronol\'ogico del \cref{rm:ch:cosmo} determina
\begin{equation}
 R_{108}=\ell_P,
 \qquad
 t_P=\frac{\ell_P}{c},
 \qquad
 T_{108}=2\pi t_P,
 \label{rm:eq:horizon-CT108-scales}
\end{equation}
y el cuanto de Planck
\begin{equation}
 E_P=\frac{\hbar}{t_P}
 =k_B\Theta_{\rm clk}\log3.
 \label{rm:eq:horizon-Planck-trit}
\end{equation}

\begin{theorem}[Celda CT108 de media acci\'on]
\label{rm:thm:horizon-half-action}
En la especializaci\'on \(R=R_{108}=\ell_P\),
\begin{align}
 \mathcal Q_H(\ell_P)
 &=\frac{c^4\ell_P}{G}=E_P,
 \label{rm:eq:horizon-QH-EP}\\
 \boxed{
 E_\Lambda=T_HS_H
 =\frac{E_P}{2}
 =\frac{k_B\Theta_{\rm clk}\log3}{2},}
 \label{rm:eq:horizon-CT108-half-energy}\\
 \boxed{E_\Lambda t_P=\frac{\hbar}{2},}
 \qquad
 \boxed{E_\Lambda T_{108}=\frac h2.}
 \label{rm:eq:horizon-half-actions}
\end{align}
\end{theorem}

\begin{proof}
La primera igualdad se obtiene sustituyendo
\(\ell_P=\sqrt{G\hbar/c^3}\) y
\(E_P=\sqrt{\hbar c^5/G}\). La
\cref{rm:eq:horizon-half-quantum-general} produce
\(E_\Lambda=E_P/2\), y
\cref{rm:eq:horizon-Planck-trit} aporta la carta informacional. Finalmente,
\[
 E_\Lambda t_P=\frac{E_Pt_P}{2}=\frac\hbar2
\]
y, como \(T_{108}=2\pi t_P\) y \(h=2\pi\hbar\),
\[
 E_\Lambda T_{108}
 =\frac{E_P}{2}(2\pi t_P)=\pi\hbar=\frac h2.
\]
\end{proof}

La primera igualdad de \cref{rm:eq:horizon-half-actions} publica media acci\'on
reducida sobre el tiempo de Planck; la segunda publica media acci\'on completa
sobre el periodo CT108. Ambas expresan una misma celda en las cartas angular
y circular.

La temperatura y la entrop\'ia quedan tambi\'en determinadas:
\begin{equation}
 k_BT_H=\frac{E_P}{2\pi},
 \qquad
 \boxed{\frac{S_H}{k_B}=\pi.}
 \label{rm:eq:horizon-CT108-TS}
\end{equation}
La temperatura convierte la frecuencia de horizonte en energ\'ia; la
entrop\'ia convierte el \'area de la esfera de Planck en informaci\'on. Su
producto recupera el mismo medio cuanto.

\subsection{Peso efectivo del horizonte y exponencial hiperb\'olica}

Definimos el peso termodin\'amico efectivo
\begin{equation}
 \Omega_{\rm eff}(R):=\exp\!\left(\frac{S_H(R)}{k_B}\right).
 \label{rm:eq:horizon-effective-weight}
\end{equation}
En el radio CT108,
\begin{equation}
 \boxed{\Omega_{\rm eff}(\ell_P)=\ee^\pi.}
 \label{rm:eq:horizon-effective-e-pi}
\end{equation}
El peso \(\Omega_{\rm eff}\) es una magnitud positiva definida por la
entrop\'ia. Su valor real expresa multiplicidad efectiva; una interpretaci\'on
como cardinal entero de microestados requiere una discretizaci\'on adicional.

El operador de incidencia areal--volum\'etrica es
\begin{equation}
 F_{\rm av}=\begin{pmatrix}0&1\\1&1\end{pmatrix},
 \qquad
 J_\varphi:=\frac{2F_{\rm av}^2-3I}{\sqrt5}.
 \label{rm:eq:horizon-J-phi}
\end{equation}
Un c\'alculo directo da
\begin{equation}
 J_\varphi^2=I,
 \qquad
 \Spec(J_\varphi)=\{+1,-1\}.
 \label{rm:eq:horizon-J-involution}
\end{equation}
El flujo hiperb\'olico en el tiempo de clausura \(\pi\) es
\begin{equation}
 \mathcal E_\pi
 :=\ee^{\pi J_\varphi}
 =\cosh\pi\,I+\sinh\pi\,J_\varphi,
 \label{rm:eq:horizon-hyperbolic-flow}
\end{equation}
y posee espectro
\begin{equation}
 \boxed{
 \Spec(\mathcal E_\pi)=\{\ee^\pi,\ee^{-\pi}\},
 \qquad
 \rho(\mathcal E_\pi)=\ee^\pi.}
 \label{rm:eq:horizon-hyperbolic-spectrum}
\end{equation}

\begin{corollary}[Confluencia escalar horizonte--Perron]
\label{rm:cor:horizon-Perron-scalar}
En el cierre CT108,
\begin{equation}
 \boxed{
 \frac{S_H}{k_B}
 =\log\Omega_{\rm eff}
 =\log\rho(\mathcal E_\pi)
 =\pi.}
 \label{rm:eq:horizon-Perron-log}
\end{equation}
\end{corollary}

\begin{proof}
Las ecuaciones
\cref{rm:eq:horizon-CT108-TS,rm:eq:horizon-effective-e-pi,rm:eq:horizon-hyperbolic-spectrum}
dan los tres miembros.
\end{proof}

La identidad escalar es exacta y conserva dos genealog\'ias distintas. El
horizonte produce \(\ee^\pi\) mediante exponenciaci\'on de la entrop\'ia
areal; la din\'amica de Perron lo produce como autovalor expansivo de una
involuci\'on hiperb\'olica. La igualdad terminal fija una condici\'on de borde
para el entrelazador y preserva la distinci\'on de dominios.

\subsection{Consistencia del transductor gravitatorio con el horizonte de Schwarzschild}

Un horizonte de Schwarzschild con el mismo radio areal satisface
\begin{equation}
 R_H=\frac{2GE_H}{c^4},
 \qquad
 E_H=\frac{c^4R_H}{2G}.
 \label{rm:eq:horizon-Schwarzschild-energy}
\end{equation}
La entrop\'ia areal coincide con \cref{rm:eq:horizon-entropy}, mientras la
temperatura es
\begin{equation}
 k_BT_{\rm Sch}=\frac{\hbar c}{4\pi R_H}
 =\frac12k_BT_H\big|_{R=R_H}.
 \label{rm:eq:horizon-Schwarzschild-temperature}
\end{equation}
En consecuencia,
\begin{equation}
 E_H=2T_{\rm Sch}S_H,
 \qquad
 T_H^{\rm dS}=2T_{\rm Sch}
 \label{rm:eq:horizon-Schwarzschild-Smarr}
\end{equation}
para un mismo radio. En \(R_H=\ell_P\),
\begin{equation}
 E_H=\frac{E_P}{2},
 \qquad
 T_{\rm Sch}S_H=\frac{E_P}{4}.
 \label{rm:eq:horizon-Schwarzschild-Planck}
\end{equation}
El contraste separa las dos geometr\'ias: comparten energ\'ia cuasilocal,
\'area y entrop\'ia, y sus gravedades superficiales producen temperaturas
en raz\'on dos.

Con \(\beta=(k_BT_{\rm Sch})^{-1}\), la acci\'on eucl\'idea satisface
\begin{equation}
 \frac{I_E}{\hbar}
 =\frac{E_H}{k_BT_{\rm Sch}}-\frac{S_H}{k_B}
 =2\pi-\pi=\pi
 \label{rm:eq:horizon-Schwarzschild-action}
\end{equation}
en el radio de Planck. El mismo \(\pi\) aparece como entrop\'ia reducida,
acci\'on eucl\'idea reducida, tiempo del flujo hiperb\'olico y logaritmo del
peso efectivo.

\subsection{Entrelazador horizonte–Perron: compatibilidad espectral entre energía de frontera y memoria de retorno}

La coincidencia
\(\Omega_{\rm eff}=\rho(\mathcal E_\pi)\) fija un dato espectral y deja
abierta la relaci\'on entre los objetos completos. Denotemos por
\(\mathcal A_H^{(m)}\) el \'algebra de observables de horizonte al nivel
non\'adico \(m\), por \(\omega_H^{(m)}\) su estado, y por
\(\mathcal P_\varphi^{(m)}\) el sistema de transferencia de Perron.

\begin{definition}[Entrelazador termodin\'amico--din\'amico]
\label{rm:def:horizon-Perron-intertwiner}
Un entrelazador admisible es una familia de aplicaciones lineales unitales
completamente positivas (UCP)
\begin{equation}
 \mathfrak I_m:\mathcal A_H^{(m)}\longrightarrow
 \mathcal A_\varphi^{(m)}
 \label{rm:eq:horizon-intertwiner}
\end{equation}
que satisface:
\begin{enumerate}[label=\roman*)]
 \item preservaci\'on del estado, adem\'as de la unidad y de la positividad
       completa ya incorporadas en la condici\'on UCP;
 \item compatibilidad con el refinamiento non\'adico;
 \item transporte de la involuci\'on interior--exterior hacia
       \(J_\varphi\mapsto-J_\varphi\);
 \item conmutaci\'on entre el semigrupo modular de horizonte y el semigrupo
       de transferencia sobre el subespacio declarado;
 \item igualdad del radio espectral en el tiempo \(\pi\), dada por
       \cref{rm:eq:horizon-Perron-log}.
\end{enumerate}
\end{definition}

La quinta condici\'on es necesaria y escalar; las cuatro primeras aportan el
contenido algebraico que una coincidencia num\'erica no contiene. Un
entrelazador construido de este modo identificar\'ia la multiplicidad
termodin\'amica efectiva con una din\'amica de expansi\'on y contracci\'on
conservando orientaci\'on, estado y refinamiento.

\begin{proofstatusbox}
Las identidades
\(E_\Lambda=E_{\rm MS}=T_HS_H=c^4R/(2G)\), la definici\'on
\(\mathcal Q_H=\mathscr F_P\,R\), la especializaci\'on CT108, las dos celdas
de media acci\'on y la igualdad
\(S_H/k_B=\log\rho(\ee^{\pi J_\varphi})=\pi\) son resultados exactos bajo
las convenciones y la realizaci\'on de Sitter declaradas. La identificaci\'on
de \(\Omega_{\rm eff}\) con un cardinal microsc\'opico y el entrelazador
\(\mathfrak I_m\) pertenecen a una interfaz f\'isica delimitada.
\end{proofstatusbox}


\begin{falsifierbox}
El cierre geom\'etrico falla si \(\Lambda_R\ne3/R^2\), si la esfera de
radio \(R\) incumple la condici\'on marginal o si la energ\'ia integrada difiere de
\(E_{\rm MS}(R)\). El cierre termodin\'amico falla si la temperatura o la
ley de \'area producen \(T_HS_H\ne c^4R/(2G)\). La especializaci\'on CT108
falla si \(R_{108}\ne\ell_P\), \(T_{108}\ne2\pi t_P\) o
\(E_P\ne k_B\Theta_{\rm clk}\log3\). La interfaz Perron queda refutada por
una obstrucci\'on de positividad, estado, orientaci\'on o refinamiento, aun
cuando el autovalor \(\ee^\pi\) coincida escalarmente.
\end{falsifierbox}
